\clearpage
\addcontentsline{toc}{chapter}{啟示錄給你的福音}
\begin{center}
\begin{tcolorbox}[
    width=0.95\textwidth,
    colback=white, %gray!5,
    colframe=gray!70,
    colbacktitle=gray,
    coltitle=white,
    title={\small{\centering 你知道啓示錄中，包含了許多神要讓你知道的福音嗎？}},
    fontupper=\small\setstretch{1.0},
    enhanced,
    boxrule=0.5pt,
    breakable,
    left=0pt,
    right=6pt,
    top=6pt,
    bottom=6pt
]

一、神是創造萬有、配得敬拜的主

你的生命是有目的。人是照著神尊貴的形象造的，並不是從猿猴演化來的。神創造萬有，祂是生命的源頭，也是全人類應當敬拜的主。
\begin{ScriptureBox}
\item[4:11] 我們的主，我們的神，你是配得榮耀、尊貴、權柄的；因為你創造了萬物，並且萬物是因你的旨意被創造而有的。
\end{ScriptureBox}

\begin{itemize}[leftmargin=2em, labelsep=0.3em]
\item[Q.] \AskItem{你相信人是照著神尊貴的形象造的嗎？}
\item[Q.] \AskItem{你知道我們需要重新認識自己的創造主嗎？}
\end{itemize}
\vspace{0.3cm}

二、沒有人可以永遠逃避神的公義。
\begin{ScriptureBox}
\item[2:23] 我是那察看人肺腑心腸的，並要照你們的行為報應你們各人。
\item[20:12-15] 我又看見死了的人，無論大小，都站在寶座前。案卷都展開了，並另有一卷展開，就是生命冊。死了的人都憑着這些案卷所記載的，照他們所行的受審判。 於是海交出其中的死人，死亡和陰間也交出其中的死人；他們都照各人所行的受審判。 死亡和陰間也被扔進火湖裏，這火湖就是第二次的死。 凡名字沒有記在生命冊上的人，就被扔進火湖裏。」
\end{ScriptureBox}

\begin{itemize}[leftmargin=2em, labelsep=0.3em]
\item[Q.] \AskItem{你相信神不只看外在行為，也鑒察人的內心嗎？}
\item[Q.] \AskItem{你知道人人都有一死，死後又有審判嗎？}
\item[Q.] \AskItem{罪不只是犯了幾個錯誤，也不只是道德上的不完全。罪更是人背離神，不願按照神的旨意生活，反而以自己為中心，追求自己的道路。}
\item[Q.] \AskItem{人可能有善良的一面，也可能做過許多好事；但我們不能因此說自己從未得罪神，也不能靠自己的善行除去罪的責任。}
\item[Q.] \AskItem{人常常以自己為中心，按照自己的意思生活。我們可能追求財富、成就、家庭與人生的滿足，卻忽略了創造我們的神，甚至把受造之物放在神的位置上。}
\end{itemize}


\vspace{0.3cm}
三、耶穌基督是神子，祂用自己的寶血洗淨我們的罪
\begin{ScriptureBox}
\item[1:5]他愛我們，用自己的血使我們脫離罪惡。
\item[5:9]因為你曾被殺、用自己的血從各族各方、各民各國中買了人來、叫他們歸於　神、
\end{ScriptureBox}
為什麼我們需要基督的寶血？因為我們都犯了罪，得罪了聖潔的神。罪不是靠行善就能抵銷的過犯。
然而，神愛罪人，沒有任憑我們留在罪中。耶穌基督為我們受死，承擔罪的刑罰，使蒙祂救贖的人得以潔淨。

\vspace{0.3cm}
四、人需要救恩
\begin{ScriptureBox}
\item[7:14]他們曾用羔羊的血把衣裳洗白淨了。
\end{ScriptureBox}
這些人能夠站在神面前，不是因為他們從未犯過罪，而是因為他們知道用羔羊的血把污穢的心洗潔淨了。

\vspace{0.3cm}
五、神呼召我們悔改、相信基督，白白領受救恩
\begin{ScriptureBox}
\item[22:17] 來！聽見的人也該說：來！口渴的人也當來；願意的，都可以白白取生命的水喝。
\item[21:4] 神要擦去他們一切的眼淚；不再有死亡，也不再有悲哀、哭號、疼痛，因為以前的事都過去了。
\end{ScriptureBox}

\begin{itemize}[leftmargin=2em, labelsep=0.3em]
\item[Q.] \AskItem{今天，你願意接受神的恩典的邀請嗎？}
\item[Q.] \AskItem{你不需要先把自己改造成一個完美的人，才可以來到祂面前。你可以承認自己的罪，悔改歸向神，相信耶穌基督為罪人受死並且復活，接受祂所賜的救恩，並開始學習跟隨祂。}
\end{itemize}

\end{tcolorbox}
\end{center}


